\section{Statement and Solution}

We work entirely within the framework established in the problem. Our goal is to give a
manifestly positive combinatorial interpretation of the universal coefficients
$c_{(a,1^b)}(n)$ when the indexing partition is a hook shape. Throughout, $a\ge 1$ and
$b\ge 0$ are integers, and $n\ge 1$.

We freely use the following ingredients, all recalled in the problem statement:
\begin{itemize}
\item the diagonal supersymmetry expansion (Proposition~\ref{prop:ds}),
\begin{equation}\label{eq:ds}
\Hilb(R_n^{(1,1)};q;u)=\sum_{\lambda\in P(1,1,n)}c_\lambda(n)\,s_\lambda(q/u);
\end{equation}
\item the Rhoades--Wilson Hilbert series formula (Theorem~\ref{thm:hilbert}),
\begin{equation}\label{eq:rw}
\Hilb(R_n^{(1,1)};q;u)=\sum_{k=0}^n [k]_q!\,S[n,k]\,u^{\,n-k};
\end{equation}
\item the hook-shape evaluation of super Schur functions (Lemma~\ref{lem:schur-1-1}),
\begin{equation}\label{eq:hookschur}
s_{(a,1^b)}(q/u)=q^a u^b+q^{a-1}u^{b+1}\quad(a\ge1,\ b\ge0),\qquad s_\varnothing(q/u)=1;
\end{equation}
\item the inversion generating function for ordered set partitions
(Proposition~\ref{prop:OSP-inv}),
\begin{equation}\label{eq:osp}
[k]_q!\,S[n,k]=\sum_{w\in\OSP(n,k)}q^{\inv(w)};
\end{equation}
\item the Sagan--Swanson involution $\phi$ on ordered set partitions, with its types I,
II, III, and the facts that $\phi$ is an involution, that merging (type I) decreases
$\inv$ by $1$ and decreases the number of blocks by $1$, that splitting (type II) increases
$\inv$ by $1$ and increases the number of blocks by $1$, and that fixed points are exactly
the type III ordered set partitions.
\end{itemize}

% --- Lemma: first (signed) formula -----------------------------------------------------
\begin{lemma}\label{lem:signed}
For $a\ge1$ and $b\ge0$,
\begin{equation}\label{eq:signedformula}
c_{(a,1^b)}(n)=\sum_{i=0}^{b}(-1)^i\,\#\{w\in\OSP(n,n-b+i)\ :\ \inv(w)=a+i\}.
\end{equation}
\end{lemma}

\begin{proof}
Equating \eqref{eq:ds} and \eqref{eq:rw}, then using \eqref{eq:hookschur} and noting that
$P(1,1,n)$ consists exactly of the empty partition together with the hooks $(a,1^b)$ of
length at most $n$, we obtain
\begin{equation}\label{eq:master}
1+\sum_{\substack{a\ge1\\ b\ge0}}c_{(a,1^b)}(n)\bigl(q^a u^b+q^{a-1}u^{b+1}\bigr)
=\sum_{d=0}^{n}[d]_q!\,S[n,d]\,u^{\,n-d}.
\end{equation}
Substituting \eqref{eq:osp} into the right-hand side of \eqref{eq:master} gives
\begin{equation}\label{eq:master2}
1+\sum_{\substack{a\ge1\\ b\ge0}}c_{(a,1^b)}(n)\bigl(q^a u^b+q^{a-1}u^{b+1}\bigr)
=\sum_{d=0}^{n}\ \sum_{i\ge0} q^{i}\,\#\{w\in\OSP(n,d):\inv(w)=i\}\,u^{\,n-d}.
\end{equation}
Both sides are polynomials in $q,u$. We extract the coefficient of $q^a u^b$ with $a\ge1$.
On the left, the monomial $q^a u^b$ arises from the term indexed by $(a,1^b)$ (through
$q^a u^b$, present since $a\ge1,b\ge0$) and, when $b\ge1$, from the term indexed by
$(a+1,1^{b-1})$ (through $q^{a}u^{b}$, since that term contributes
$q^{a+1}u^{b-1}+q^{a}u^{b}$). The constant term $1$ contributes nothing for $a\ge1$. On the
right, the coefficient of $u^b$ corresponds to $d=n-b$, and then the coefficient of $q^a$
is the number of $w\in\OSP(n,n-b)$ with $\inv(w)=a$. Hence:
\begin{align}
c_{(a,1^b)}(n)+c_{(a+1,1^{b-1})}(n)&=\#\{w\in\OSP(n,n-b):\inv(w)=a\}\qquad(b\ge1),
\label{eq:rec1}\\
c_{(a)}(n)&=\#\{w\in\OSP(n,n):\inv(w)=a\}\qquad(b=0).
\label{eq:rec0}
\end{align}
Set
\[
N(a,b):=\#\{w\in\OSP(n,n-b):\inv(w)=a\}.
\]
Then \eqref{eq:rec1} reads $c_{(a,1^b)}(n)=N(a,b)-c_{(a+1,1^{b-1})}(n)$ for $b\ge1$, and
\eqref{eq:rec0} reads $c_{(a)}(n)=N(a,0)$. Iterating the recurrence,
\begin{align*}
c_{(a,1^b)}(n)
&=N(a,b)-c_{(a+1,1^{b-1})}(n)\\
&=N(a,b)-N(a+1,b-1)+c_{(a+2,1^{b-2})}(n)\\
&=\cdots=\sum_{i=0}^{b}(-1)^i\,N(a+i,\,b-i),
\end{align*}
where the final term, at $i=b$, is $(-1)^b c_{(a+b)}(n)=(-1)^b N(a+b,0)$, consistent with
the pattern. Since $N(a+i,b-i)=\#\{w\in\OSP(n,n-(b-i)):\inv(w)=a+i\}
=\#\{w\in\OSP(n,n-b+i):\inv(w)=a+i\}$, this is exactly \eqref{eq:signedformula}.
\end{proof}

% --- Cancellation set and partition --------------------------------------------------
We now convert the alternating sum \eqref{eq:signedformula} into a positive sum via the
Sagan--Swanson sign-reversing involution. Fix integers $b\ge0$ and $\ell\ge1$ with
$0\le b\le \lfloor (\ell-1)/2\rfloor$, and recall the set
\begin{equation}\label{eq:Sdef}
S_{n,\ell,b}=\bigcup_{c=0}^{b}\{w\in\OSP(n,n-c):\inv(w)=\ell-b-c\}.
\end{equation}
The union in \eqref{eq:Sdef} is disjoint: two members of the union with indices
$c\ne c'$ lie in $\OSP(n,n-c)$ and $\OSP(n,n-c')$ respectively, which are disjoint sets
since the number of blocks differs. Define the signed weight on ordered set partitions by
$\sgn(w)=(-1)^{\,n-\blocks(w)}$. For $w\in\OSP(n,n-c)$ we have $\blocks(w)=n-c$, so
$\sgn(w)=(-1)^c$.

Comparing \eqref{eq:Sdef} (with $\ell=a+2b$) to \eqref{eq:signedformula}: setting $i=c$ and
$\ell=a+2b$, the set $\{w\in\OSP(n,n-b+i):\inv(w)=a+i\}$ has $\blocks(w)=n-(b-i)=n-c'$ for
$c'=b-i$, and $\inv(w)=a+i=\ell-b-c'$. Re-indexing $c=b-i$ runs $c$ over $0,\dots,b$ as $i$
does, and the sign $(-1)^i=(-1)^{b-c}=(-1)^b(-1)^c$. Therefore
\begin{equation}\label{eq:cassigned}
c_{(a,1^b)}(n)=\sum_{i=0}^{b}(-1)^i N(a+i,b-i)
=(-1)^b\sum_{w\in S_{n,a+2b,b}}\sgn(w).
\end{equation}

% --- Lemma describing how phi acts on S -----------------------------------------------
\begin{lemma}\label{lem:phiaction}
Fix $\ell\ge1$ and $0\le b\le\lfloor(\ell-1)/2\rfloor$. Let $w\in S_{n,\ell,b}$, say
$w\in\OSP(n,n-c)$ with $\inv(w)=\ell-b-c$, where $0\le c\le b$. Then:
\begin{enumerate}[(i)]
\item $w$ cannot be of type III;
\item if $w$ is of type I and $c=b$ (equivalently $\inv(w)=\ell-2b$), then
$\phi(w)\notin S_{n,\ell,b}$;
\item if $w$ is of type I and $c<b$, then $\phi(w)\in S_{n,\ell,b}$;
\item if $w$ is of type II, then $\phi(w)\in S_{n,\ell,b}$.
\end{enumerate}
\end{lemma}

\begin{proof}
Note first that $w\in S_{n,\ell,b}$ implies, as $c$ ranges over $0,\dots,b$ and
$\inv(w)=\ell-b-c$, the bounds
\begin{equation}\label{eq:invbounds}
\ell-2b\le \inv(w)\le \ell-b .
\end{equation}

(i) If $w$ were of type III, then $w$ is a fixed point of $\phi$, i.e.\ no $M$ exists at
which $w$ is splittable or mergeable. By the analysis of Sagan--Swanson, a type III
ordered set partition must have all blocks of size $1$ (otherwise it would be splittable),
hence $\blocks(w)=n$, forcing $c=0$. Then $\inv(w)=\ell-b$, and being of type III with all
singleton blocks and no mergeable index means the maxima of the singletons are increasing,
so $\inv(w)=0$, giving $\ell=b$. But then $b=\ell\ge1$, whence
$b>\lfloor(\ell-1)/2\rfloor$, contradicting $b\le\lfloor(\ell-1)/2\rfloor$. Thus $w$ is not
of type III.

(ii) Here $c=b$, so $\inv(w)=\ell-2b$ and $w\in\OSP(n,n-b)$. Since $w$ is of type I,
applying $\phi$ merges two blocks: $\inv$ decreases by $1$ and the number of blocks
decreases by $1$. Thus $\phi(w)\in\OSP(n,n-b-1)$ with $\inv(\phi(w))=\ell-2b-1$. For
$\phi(w)$ to lie in $S_{n,\ell,b}$ we would need $\phi(w)\in\OSP(n,n-c')$ with
$0\le c'\le b$ and $\inv(\phi(w))=\ell-b-c'$. The block count forces $c'=b+1$, which is
$>b$, so $\phi(w)\notin S_{n,\ell,b}$.

(iii) Here $c<b$ and $w$ is type I. After merging, $\phi(w)\in\OSP(n,n-(c+1))$ with
$\inv(\phi(w))=(\ell-b-c)-1=\ell-b-(c+1)$. Since $0\le c+1\le b$, the element $\phi(w)$ lies
in the $c+1$ summand of \eqref{eq:Sdef}, i.e.\ $\phi(w)\in S_{n,\ell,b}$.

(iv) Here $w$ is type II, so applying $\phi$ splits a block: $\inv$ increases by $1$ and the
block count increases by $1$. We have $w\in\OSP(n,n-c)$; since splitting is possible some
block has size $>1$, so $\blocks(w)<n$, i.e.\ $c\ge1$. After splitting,
$\phi(w)\in\OSP(n,n-(c-1))$ with $\inv(\phi(w))=(\ell-b-c)+1=\ell-b-(c-1)$. Since
$0\le c-1\le b$, $\phi(w)$ lies in the $c-1$ summand of \eqref{eq:Sdef}, i.e.\
$\phi(w)\in S_{n,\ell,b}$.
\end{proof}

% --- Main theorem ---------------------------------------------------------------------
\begin{theorem}\label{thm:main}
For $a\ge1$ and $b\ge0$,
\begin{equation}\label{eq:posformula}
c_{(a,1^b)}(n)=\#\{w\in\OSP(n,n-b)\ :\ \inv(w)=a\ \text{and}\ w\ \text{is of type I}\}.
\end{equation}
In particular $c_{(a,1^b)}(n)\ge0$, recovering its known nonnegativity in a manifestly
positive form.
\end{theorem}

\begin{proof}
Take $\ell=a+2b$. Then $a\ge1$ gives $\ell\ge1$ and $\ell-1=a+2b-1\ge 2b$, so
$b\le\lfloor(\ell-1)/2\rfloor$; thus Lemma~\ref{lem:phiaction} applies. By
\eqref{eq:cassigned},
\begin{equation}\label{eq:signedsum}
c_{(a,1^b)}(n)=(-1)^b\sum_{w\in S_{n,a+2b,b}}\sgn(w).
\end{equation}

We now show $\phi$ restricts to a sign-reversing, almost-fixed-point-free involution on
$S_{n,a+2b,b}$, whose only ``defect'' is a single class of elements on which $\phi$ leaves
$S_{n,a+2b,b}$. Concretely, define
\begin{equation}\label{eq:Tdef}
T_{n,a+2b,b}:=\{w\in\OSP(n,n-b)\ :\ \inv(w)=a\ \text{and}\ w\ \text{is of type I}\}.
\end{equation}
Note $T_{n,a+2b,b}\subseteq S_{n,a+2b,b}$: an element $w$ of $T_{n,a+2b,b}$ lies in
$\OSP(n,n-b)$ (so $c=b$) and has $\inv(w)=a=\ell-2b=\ell-b-c$, matching the $c=b$ summand of
\eqref{eq:Sdef}.

We partition $S_{n,a+2b,b}$ according to Lemma~\ref{lem:phiaction}:
\begin{itemize}
\item By Lemma~\ref{lem:phiaction}(i), no element of $S_{n,a+2b,b}$ is of type III, hence no
element of $S_{n,a+2b,b}$ is a fixed point of $\phi$. Thus every element is of type I or
type II.
\item By Lemma~\ref{lem:phiaction}(iv), if $w\in S_{n,a+2b,b}$ is of type II then
$\phi(w)\in S_{n,a+2b,b}$.
\item By Lemma~\ref{lem:phiaction}(iii), if $w\in S_{n,a+2b,b}$ is of type I and lies in a
summand with $c<b$, then $\phi(w)\in S_{n,a+2b,b}$.
\item By Lemma~\ref{lem:phiaction}(ii), if $w\in S_{n,a+2b,b}$ is of type I and lies in the
summand with $c=b$ (equivalently $w\in\OSP(n,n-b)$, $\inv(w)=a$), then
$\phi(w)\notin S_{n,a+2b,b}$. These $w$ are exactly the elements of $T_{n,a+2b,b}$.
\end{itemize}

Let $U:=S_{n,a+2b,b}\setminus T_{n,a+2b,b}$. We claim $\phi$ maps $U$ to $U$ and is an
involution there. Indeed:
\begin{itemize}
\item If $w\in U$ is of type II, then $\phi(w)\in S_{n,a+2b,b}$ by the above. Moreover
$\phi(w)$ is of type I (since $\phi$ is an involution interchanging types I and II among
non-fixed points), and $\phi(w)$ lies in $\OSP(n,n-(c-1))$ with $c\ge1$; if
$\phi(w)\in T_{n,a+2b,b}$ then $\phi(w)\in\OSP(n,n-b)$ would force $c-1=b$, i.e.\ $c=b+1$,
impossible since $c\le b$. Hence $\phi(w)\in U$.
\item If $w\in U$ is of type I, then by definition of $U$ it is not in $T_{n,a+2b,b}$, so it
lies in a summand with $c<b$; by Lemma~\ref{lem:phiaction}(iii), $\phi(w)\in S_{n,a+2b,b}$,
and $\phi(w)$ is of type II, hence $\phi(w)\notin T_{n,a+2b,b}$ (whose elements are type I),
so $\phi(w)\in U$.
\end{itemize}
Since $\phi$ is an involution on all ordered set partitions and we have shown it preserves
$U$, it is an involution on $U$. It has no fixed points on $U$ (no type III elements occur
in $S_{n,a+2b,b}$ at all). Finally, $\phi$ is sign-reversing for $\sgn(w)=(-1)^{n-\blocks(w)}$:
merging or splitting changes $\blocks(w)$ by exactly $1$, so $\sgn(\phi(w))=-\sgn(w)$ for
all non-fixed points. Therefore the elements of $U$ cancel in pairs in the signed sum:
\[
\sum_{w\in U}\sgn(w)=0 .
\]

Consequently, from \eqref{eq:signedsum},
\[
c_{(a,1^b)}(n)=(-1)^b\sum_{w\in S_{n,a+2b,b}}\sgn(w)
=(-1)^b\Bigl(\sum_{w\in U}\sgn(w)+\sum_{w\in T_{n,a+2b,b}}\sgn(w)\Bigr)
=(-1)^b\sum_{w\in T_{n,a+2b,b}}\sgn(w).
\]
Every $w\in T_{n,a+2b,b}$ lies in $\OSP(n,n-b)$, so $\blocks(w)=n-b$ and
$\sgn(w)=(-1)^{n-(n-b)}=(-1)^b$. Hence $(-1)^b\sgn(w)=(-1)^{2b}=1$ for each such $w$, and
\[
c_{(a,1^b)}(n)=\sum_{w\in T_{n,a+2b,b}}1=\#T_{n,a+2b,b}
=\#\{w\in\OSP(n,n-b):\inv(w)=a,\ w\ \text{type I}\},
\]
which is exactly \eqref{eq:posformula}. As this is a cardinality, it is nonnegative.
\end{proof}

\begin{remark}
For $b=0$ the formula \eqref{eq:posformula} reads
$c_{(a)}(n)=\#\{w\in\OSP(n,n):\inv(w)=a,\ w\ \text{type I}\}$. Here every block is a
singleton; an ordered set partition into singletons is determined by a permutation, and it
is of type I precisely when it is mergeable, i.e.\ when it is not increasing-from-the-last
in the relevant sense. This matches \eqref{eq:rec0}: indeed for these, all elements of
$\OSP(n,n)$ with $\inv=a\ge1$ are of type I except possibly the single increasing one, which
has $\inv=0$, so for $a\ge1$ the type I restriction is automatic, recovering
$c_{(a)}(n)=\#\{w\in\OSP(n,n):\inv(w)=a\}$, consistent with \eqref{eq:rec0}. The table for
$n=4$ in the problem statement illustrates \eqref{eq:posformula} in all cases.
\end{remark}

\begin{corollary}\label{cor:general-m}
For any $m,n\ge1$ and any hook shape $\lambda=(a,1^b)$ with $a\ge1$, $b\ge0$ and
$\ell(\lambda)=b+1\le n$, the coefficient $c_\lambda(n)$ appearing in
\[
\Hilb(R_n^{(m)};q_1,\dots,q_m)=\sum_{\lambda}c_\lambda(n)\,s_\lambda(q_1,\dots,q_m)
\]
is given by the manifestly positive combinatorial formula \eqref{eq:posformula}.
\end{corollary}

\begin{proof}
By the cited results of Bergeron and Lentfer, the coefficients $c_\lambda(n)$ are
independent of $m$, and for hook shapes they coincide with the coefficients
$c_\lambda(n)$ of Proposition~\ref{prop:ds} attached to the superspace coinvariant ring
$R_n^{(1,1)}$, since $P(1,1,n)$ is precisely the set of hooks of length at most $n$. The
formula \eqref{eq:posformula} of Theorem~\ref{thm:main} therefore computes them.
\end{proof}

This completes the combinatorial interpretation: for a hook shape $\lambda=(a,1^b)$,
$c_\lambda(n)$ is the number of ordered set partitions of $\{1,\dots,n\}$ into $n-b$ blocks
with exactly $a$ inversions which are of type I (mergeable at their largest splittable or
mergeable maximum) in the sense of Sagan--Swanson.